\begin{lemma}\label{centralizers}
Let $M=\operatorname{PSL}_3(7):\langle\tau\rangle$, where
$\tau$ is inverse transpose. An inner involution has centralizer of
order 1344 in $M$. The centralizer of an outer involution has no
element of order 16.
\end{lemma}
\begin{proof}
An inner involution lifts, after multiplication by a scalar cube
root of unity, to an involution of $\operatorname{SL}_3(7)$. It is
conjugate to $\operatorname{diag}(-1,-1,1)$. A projective centralizer
cannot send this matrix to a nontrivial scalar multiple: their traces
would differ. Its centralizer in $\operatorname{SL}_3(7)$ is the block
group $\operatorname{GL}_2(7)$, of order 2016. Quotienting by the
scalar centre of order three gives 672. Inverse transpose centralizes
the diagonal involution, so the centralizer in $M$ has order 1344.

Write an outer involution as $z=[A]\tau$ with $A\in
\operatorname{SL}_3(7)$. The involution condition gives
$AA^{-T}=\zeta I$ with $\zeta^3=1$, or $A=\zeta A^T$.
Transposing again gives $\zeta^2=1$, hence $\zeta=1$ and $A$ is
symmetric. The inner centralizer is the projective similitude group
of this nondegenerate three-dimensional symmetric form. Every
similitude multiplier $\mu$ is a square, since its determinant
satisfies $(\det h)^2=\mu^3$. Thus each projective similitude has a
unique representative in the special orthogonal group: first rescale
to an isometry, then choose the sign making its determinant one.
These representatives belong to $\operatorname{SL}_3(7)$, so no
additional restriction arises from passing to $\operatorname{PSL}_3(7)$.

A semisimple element of 2-power order in this odd-dimensional special
orthogonal group has eigenvalues $1,\lambda,\lambda^{-1}$. Frobenius
either fixes the two nontrivial eigenvalues or exchanges them. Thus
their orders divide 6 or 8, respectively, and the element's 2-power
order divides 8. Finally $C_M(z)=C_{\operatorname{PSL}_3(7)}(z)
\times\langle z\rangle$, so adjoining the outer involution still
introduces no element of order 16.
\end{proof}
